\documentclass{article}
\usepackage{amsmath,amssymb,hyperref}
\begin{document}
\section*{Continuity and strict monotonicity of Nyman--Beurling distances}
\subsection{Definitions and mathematical statement}
Work in the Hilbert space
\[
 H=L^2((0,\infty),t^{-2}\,dt),\qquad
 \|f\|_H^2=\int_0^\infty |f(t)|^2\,\frac{dt}{t^2}.
\]
Let $\chi$ be the indicator of $[1,\infty)$ and put
$e_\alpha(t)=\{t/\alpha\}$, where braces denote fractional part.
For real $\lambda\geq1$, define
\[
 D(\lambda)=\inf\{\|\chi-v\|_H:
       v\in\operatorname{span}\{e_\alpha:1\leq\alpha\leq\lambda\}\}.
\]
The span uses finite linear combinations over $\mathbb R$; allowing complex
coefficients or taking its closure gives the same distance.

Is $D$ continuous on $[1,\infty)$, including right-continuity at $1$?
Is it strictly decreasing there, so that
$D(\lambda_1)>D(\lambda_2)$ whenever $1\leq\lambda_1<\lambda_2$?

\subsection{Short English statement}
As the permitted range of fractional-part dilations expands, is the distance
to the indicator function continuous and strictly decreasing?

\subsection{Sources}
M. Balazard's second and third questions in the problem session of
\emph{Dirichlet Series and Function Theory in Polydiscs}, Oberwolfach
Report 6/2014, p.\ 390, section 8; reporter O. F. Brevig,
\url{https://doi.org/10.4171/owr/2014/06}.
This is an independently written formulation of the two questions.
\end{document}
